% Generated draft: requires mathematical and source review.
\documentclass[11pt]{article}
\usepackage[margin=25mm]{geometry}
\usepackage{fontspec}
\setmainfont{Latin Modern Roman}
\usepackage{amsmath,amssymb,mathtools,unicode-math}
\setmathfont{Latin Modern Math}
\usepackage{enumitem,needspace,xurl,hyperref}
\setlength{\parindent}{0pt}
\setlength{\parskip}{5pt}
\setlength{\emergencystretch}{3em}
\newcommand{\N}{\mathbb{N}}
\newcommand{\Z}{\mathbb{Z}}
\newcommand{\Q}{\mathbb{Q}}
\newcommand{\R}{\mathbb{R}}
\newcommand{\C}{\mathbb{C}}
\newcommand{\F}{\mathbb{F}}
\newcommand{\A}{\mathbb{A}}
\newcommand{\PP}{\mathbb{P}}
\newcommand{\E}{\mathbb{E}}
\newcommand{\eps}{\varepsilon}
\newcommand{\dd}{\,\mathrm{d}}
\begin{document}
\section*{UnsolvedMath ID: 30000305}
\textbf{Source problem code:} catalogue entry
\subsection{Definitions and mathematical statement}
\textbf{Research problem.} Determine how multiple mergers of ancestral lineages affect the full site-frequency spectrum, including the number of mutations affecting exactly \(k\) sampled lineages for each \(k\). For a sample of \(n\) labeled genomes, let \(\xi_k^{(n)}\) count mutations carried by exactly \(k\) sampled descendants. In the usual unfolded segregating-site convention the spectrum is \((\xi_1^{(n)},\ldots,\xi_{n-1}^{(n)})\); mutations carried by all \(n\) descendants require a separate fixed-site or ancestral-branch convention. The count \(k\) refers to sampled descendants of a mutated branch, not to the number of ancestral blocks participating in one merger event.\par A conditional coalescent formulation uses a partition process \(\Pi_t\) of \(\{1,\ldots,n\}\), whose blocks record common ancestors. A multiple merger combines more than two blocks, while simultaneous distinct mergers require additional specification. If \(T\) is the most recent common-ancestor time, define\[L_k^{(n)}=\int_0^T\#\{B\in\Pi_t:|B|=k\}\,dt.\] Under a specified infinite-sites model with independent Poisson mutations at rate \(\rho\) per unit branch length, the conditional spectrum consists of independent Poisson counts with means \(\rho L_k^{(n)}\). In particular, when these expectations are finite, \(\mathbb E\xi_k^{(n)}=\rho\,\mathbb E L_k^{(n)}\). These conventions explain one precise connection between merger mechanisms and frequency spectra; they are not supplied assumptions of every model in the question. A full joint distribution requires the joint branch-length law, not merely its expectations.\par The source does not specify the merger rates or their \(\Lambda\)- or \(\Xi\)-coalescent class, the mutation mechanism, population-size time scaling, sample size regime, unfolded versus folded data, or whether \emph{full} means all expected entries or the joint law. Nor is a comparison with binary mergers normalized. Consequently no universal direction of change or single numerical formula is asserted for every merger mechanism.\par The supplied triage checked 2026-08-21 labels the problem partially solved, reporting computed site-frequency spectra in several multiple-merger models while describing a general treatment for arbitrary mechanisms as active. No specific models or theorem hypotheses are supplied, so this dated assessment requires independent verification.
\subsection{Short English statement}
Describe the mutation-frequency counts across all descendant sample sizes under multiple-merger genealogies. The mechanism and mutation model must be fixed, and expectations should be distinguished from the full joint spectrum. A mutation's frequency counts sampled descendants, not the number of merging ancestors.
\subsection{Sources}
\begin{itemize}\item Local UnsolvedMath registry, permanent ID 30000305, code catalogue entry; \emph{Site-Frequency Spectra under Multiple Ancestral Mergers}, Oberwolfach Reports Open Problems.\item \emph{Mathematical Population Genetics (2006)}, Oberwolfach Reports, DOI 10.4171/owr/2005/40, pp. 2241--2304. \url{https://doi.org/10.4171/owr/2005/40}. Proposed year and year tag: 2005; no proposer named.\item Supplied literature label \emph{original Oberwolfach report}: \url{https://doi.org/10.4171/owr/2005/40}. No separate references identifying the computed special models are supplied.\item Literature triage checked 2026-08-21, status \texttt{partially\_solved}; record updated 2026-09-16. The \texttt{corrected\_verified} field concerns notation for the lineage-count variable. Verify dated claims independently.\end{itemize}
\end{document}
